\section{Image-Based Plant Phenotyping}

\subsection{Image-Based Plant Phenotyping Versus Traditional Methods} 
\paragraph{}
Image-based plant phenotyping methods have several advantages over traditional methods. Traditional plant phenotyping methods are destructive, prone to human error, time-consuming, and challenging to perform from start to end of a plant's ontogenesis \cite{arnalbarbedoDigitalImageProcessing2013, zhangHighthroughputPhenotypingPlant2023, laxmanNoninvasiveQuantificationTomato2018}. Recent advances in image sensing, processing, and data management have produced high-throughput image-based plant phenotyping systems, which are amenable to automation and can analyse large samples of data precisely in a very short period. These image processing solutions reduce the need for potentially erroneous, time-consuming, expensive human labour in plant phenotyping \cite{daschoudhuryLeveragingImageAnalysis2019}. Despite these advantages, several aspects of image processing, which determine the overall cost, efficiency, and applicability of image processing in plant phenotyping contexts, are important. The most important of these aspects are image acquisition and processing tasks.
 
 \paragraph{}
 These Modern high-throughput image-based plant phenotyping systems are examples of cyberinfrastructure (CI) systems \cite{abebeImageBasedHighThroughputPhenotyping2023}, which are research environments that support data acquisition, storage, management and computing and processing services distributed over some network \cite{nabwireReviewApplicationArtificial2021}. For image-based plant phenotyping, this means that image acquisition and processing techniques are vital.
  
\subsection{Image Acquistion Techniques}

\subsubsection{The importance of Image Acquistion in Plant Phenotyping}
\paragraph{}
Image processing techniques in plant phenotyping cannot be applied until suitable images are acquired, and the choice of imaging technique for data acquisition affects the types of phenotypical parameters that can be measured. Image sensor advancements have provided the foundations for detecting various information from plants using electromagnetic radiation. Plant tissues interact with this radiation by absorbing, reflecting, emitting, transmitting, and fluorescing light. Different ranges of wavelengths correspond to different biochemical or physiological aspects of plant tissue, which are not visible to the naked eye \cite{liReviewImagingTechniques2014}. Therefore, the appropriate selection of image acquisition techniques is crucial for analysing different types of plant traits.

\subsubsection{Common Image Sensor Techniques In Plant Phenotyping}
\paragraph{}
Commonly used image acquisition techniques include visible, fluorescence, thermal, tomographic, and hyperspectral imaging \cite{qiuSensorsMeasuringPlant2018, liReviewImagingTechniques2014}. The most popular techniques are visible and hyperspectral imaging due to their affordability, robustness and flexibility in measuring a wide range of phenotypic parameters. Visible imaging is an affordable option that utilises high-resolution RGB cameras. These devices are sensitive to the visible spectral range and have been used to measure various phenotypic parameters like shoot biomass, yield and leaf morphology, and salinity tolerance \cite{golzarianAccurateInferenceShoot2011, hoyos-villegasGroundBasedDigitalImaging2014}. Hyperspectral imaging is another widely used in plant phenotyping. In this technique, sensors capture multiple spectral ranges, allowing for detailed measurements of various phenotypic parameters. Notable applications of hyperspectral imaging include estimating nutrient content and detecting plant responses to abiotic or biotic stressors \cite{saricApplicationsHyperspectralImaging2022}. The other imaging techniques are not as widely used as visible or hyperspectral imaging; however, they are also effective in measuring plant phenotypic parameters. For instance, thermal imaging in plant water stress relations \cite{liReviewImagingTechniques2014, ishimweApplicationsThermalImaging2014}, fluorescence imaging in plant pathogenic symptom monitoring \cite{lichtenthalerRoleChlorophyllFluorescence1988, swarbrickMetabolicConsequencesSusceptibility2006}, and tomography in estimating plant organ water content \cite{windtMRILongdistanceWater2006}.

\subsection{Image Processing Techniques}

\subsubsection{The importance of Image Segmentation in Plant Phenotyping}
\paragraph{}
There is a lack of generalised or standardised image-processing techniques in plant phenotyping; hence, bespoke image-processing pipelines are designed for specific plant species or genotypes, such as Arabidopsis \cite{liReviewImagingTechniques2014, arvidssonGrowthPhenotypingPipeline2011a}. However, segmentation is a key image-processing task common to many image-based plant phenotyping solutions. Moreover, the validity and precision of the derived phenotypical measurements and analysis depend on the image segmentation method selected \cite{choudhurySegmentationTechniquesChallenges2020, muller-linowPlantScreenMobile2019}, because the accuracy of more complex tasks, such as object detection, 3D reconstruction, and classification, depends on segmentation accuracy. Therefore, image segmentation methods are integral to image processing in plant phenotyping.

\subsubsection{Commonly Used Image Segmentation Techniques}
\paragraph{}
In plant phenotyping, selecting the appropriate segmentation technique is crucial when considering the type of phenotypic parameter to be measured. Segmentation methods are typically classifiable as traditional or learning-based \cite{choudhurySegmentationTechniquesChallenges2020}. Traditional segmentation methods typically target holistic phenotypical traits like plant height and include frame differencing \cite{daschoudhuryHolisticComponentPlant2018}, colour-based \cite{hamudaSurveyImageProcessing2016a}, edge detection \cite{wangImageSegmentationOverlapping2018}, spectral band difference-based \cite{gampaDatadrivenApproachDetecting2019}, shape modelling-based \cite{chenLeafSegmentationFunctional2019}, and graph-based segmentation \cite{praveenkumarImageBasedLeaf2019}. In contrast to traditional techniques, learning-based methods, which utilise machine learning or deep learning \cite{papeUtilizingMachineLearning2015, abdallaFinetuningConvolutionalNeural2019}, are more suitable for targeting component-based traits of specific organs or parts of a plant \cite{choudhurySegmentationTechniquesChallenges2020}.